\section{Persistence under a small attractive perturbation}
\label{sec:convex}

The initial transform has three separated groups of supporting gradients.
We now show how to retain this structure when an ideal pressure is replaced
by a sufficiently accurate physical pressure and a smaller attractive
quadratic is added. The approximation error must be small compared with
the strength of that quadratic. The proof examines both operations at the
same scale; ordinary uniform convergence alone would not retain a corner.

Throughout this section, $t$ is fixed as in
Proposition~\ref{prop:initial-triple}. A physical pressure means a function
$Q=Q_\varphi$ with $\varphi$ in the class~\eqref{eq:physical-class}.
We use the following conclusions of
Propositions~\ref{prop:thermodynamics} and~\ref{prop:transform}:
$Q$ is finite, jointly convex, and
nondecreasing in $h$; its partial conjugate $H_Q(\beta,d)$ is finite and
concave in $\beta$ for every $d\geq0$; and the suprema defining
$\T_aQ$ are attained. On every compact subset of $\Omega$, their optimizing
densities are bounded uniformly over this physical class and over
$0\leq a\leq\min(1,A+1)$. These uniform bounds use the fixed value of $t$.
They are the only model-dependent bounds needed below.

For a nonempty compact convex set $K\subset\R^2$, define its support
function and its exposed faces by
\begin{equation}
 h_K(\xi)=\max_{v\in K}\xi\cdot v,
 \qquad
 F_\xi=\{v\in K:\xi\cdot v=h_K(\xi)\}.
 \label{eq:support-face}
\end{equation}
The face $F_0$ is the whole set $K$. Two facts about finite convex
functions will be used repeatedly. Locally uniform convergence bounds
subgradients on smaller compact subsets, and limits of subgradients at
converging points belong to the subdifferential of the limiting function.
Both statements follow directly from supporting inequalities and the
bound of a slope by the oscillation on a surrounding ball. In particular,
$\partial h_K(\xi)=F_\xi$. We also use the reachable-gradient identity
$\partial F(\lambda)=\conv G_F(\lambda)$ recalled in the preceding section.

\begin{lemma}[Persistence of three gradient groups]
\label{lem:persistence}
Let $W\subset\Omega$ be compact, let $P$ be finite and convex on $\Omega$,
and let $\lambda_0=(\beta_0,h_0)\in\operatorname{int}W$. Suppose that
\begin{equation}
 \partial P(\lambda_0)=K=\conv(K_1\cup K_2\cup K_3),
 \qquad B(g,\gamma)\subset K,
 \label{eq:persistence-input}
\end{equation}
where $g\in\R^2$, $\gamma>0$, and the $K_i$ are nonempty compact convex
sets whose projections onto the density coordinate are nonnegative and
strictly ordered. Let $\zeta>0$ be such that the sets
$K_i+\overline B(0,\zeta)$ still have disjoint density projections, and
$g$ lies outside the convex hull of any two of these enlarged sets.

There exist $S<\infty$ and $a_*>0$ with the following property. For every
$0<a<a_*$ and every physical pressure $Q$ satisfying
\begin{equation}
 \|Q-P\|_W\leq a^2,
 \label{eq:persistence-error}
\end{equation}
there is a point $\lambda'\in\operatorname{int}W$ with
$|\lambda'-\lambda_0|\leq Sa$ such that
\begin{equation}
 g\in\partial(\T_aQ)(\lambda')
   =\conv(J_1\cup J_2\cup J_3).
 \label{eq:persistence-output}
\end{equation}
Here each $J_i$ is nonempty, compact, and convex, has nonnegative density
projection, and satisfies $J_i\subset K_i+\overline B(0,\zeta)$.
The constants may depend on the displayed data and the fixed physical
class, but not on $Q$. We may require $a_*\leq\min(1,A+1)$.
\end{lemma}

\begin{proof}
We prove the conclusion along an arbitrary sequence $a_m\downarrow0$
and physical pressures $Q_m$ with
$\|Q_m-P\|_W\leq a_m^2$. The radius $S$ will be chosen from the fixed
limiting data. A contradiction argument will then supply one threshold
$a_*$ valid for all admissible $Q$.

\smallskip
\noindent\textit{The geometric effect of the quadratic.}
Write $\xi=(x,y)$ and $\pi_2(s,d)=d$. Define the finite convex function
\begin{equation}
 r(x,y)=\max_{(s,d)\in K}
       \left(xs+yd+\frac{\beta_0}{2}d^2\right)
 \label{eq:scaled-transform-limit-definition}
\end{equation}

This is the function that will arise when the attractive mass and the
distance to $\lambda_0$ tend to zero at the same scale. Before proving that
limit, we identify which points of $K$ can attain its maximum.
Every maximizer of
$\xi\cdot(s,d)+\beta_0d^2/2$ over $K$ belongs to one of the $K_i$.
Indeed, write such a maximizer as a finite convex combination
$v=\sum_j\theta_jv_j$, with $v_j=(s_j,d_j)\in\bigcup_iK_i$ and
positive weights. If two of the $d_j$ differ, strict convexity of the
quadratic gives
\[
 \xi\cdot v+\frac{\beta_0}{2}(\pi_2v)^2
   <\sum_j\theta_j\left(\xi\cdot v_j+
                                      \frac{\beta_0}{2}d_j^2\right)
   \leq r(\xi),
\]
contradicting maximality. All $d_j$ must therefore coincide. The disjoint
density projections put all $v_j$ in one $K_i$, and convexity of that
$K_i$ gives $v\in K_i$. 
Thus the positive quadratic rules out a maximizing mixture across the three
density groups. To use this observation for physical pressures, we must show
that limits of their transformed gradients are themselves maximizing points
of $K$; membership in the convex hull of such points would not suffice.

\smallskip
\noindent\textit{Step 1: the pressure at scale $a_m$.}
On the expanding open set
$\Omega_m=\{\xi\in\R^2:\lambda_0+a_m\xi\in\Omega\}$, put
\begin{equation}
 q_m(\xi)=
 \frac{Q_m(\lambda_0+a_m\xi)-P(\lambda_0)}{a_m}.
 \label{eq:scaled-pressure}
\end{equation}
Every fixed compact set of scaled coordinates is mapped into
$\operatorname{int}W$ for all sufficiently large $m$. On that compact
set, the difference between $q_m$ and the corresponding rescaling of $P$
is at most $a_m$. The rescaling of $P$ converges to its directional
derivative at $\lambda_0$, which is $h_K$. This convergence is uniform
on compact sets: the rescaled functions have a common local Lipschitz
bound, and their pointwise limits are the directional derivatives.
Consequently
\begin{equation}
 q_m\longrightarrow h_K
 \quad\text{locally uniformly on }\R^2.
 \label{eq:scaled-pressure-limit}
\end{equation}

The scaling preserves subgradient vectors:
\[
 \partial q_m(\xi)
   =\partial Q_m(\lambda_0+a_m\xi).
\]
It follows that, whenever $\xi_m\to\xi$ and
$v_m\in\partial Q_m(\lambda_0+a_m\xi_m)$, the vectors $v_m$ are bounded
and every limit belongs to $F_\xi$. To see the last assertion directly,
pass the supporting inequalities of $q_m$ to the limit. They give
$v\in\partial h_K(\xi)=F_\xi$. This observation applies to faces of any
dimension; no uniqueness of the supporting vector is assumed.

\smallskip
\noindent\textit{Step 2: the transform at the same scale.}
Put
\[
 B_m(\xi)=
 \frac{\T_{a_m}Q_m(\lambda_0+a_m\xi)-P(\lambda_0)}{a_m}.
\]
We will prove that
\begin{equation}
 B_m\longrightarrow r
 \quad\text{locally uniformly on }\R^2.
 \label{eq:scaled-transform-limit}
\end{equation}
To keep track of maximizing densities as well as function values, we use the slices of $K$.
For fixed $x\in\R$, define
\[
 k_x(d)=\max\{xs:(s,d)\in K\},
\]
with value $-\infty$ when $d$ is outside the density projection of $K$.
The function $k_x$ is upper semicontinuous and concave on its compact
interval of finite values. Partial conjugacy of the support function gives
\begin{equation}
 k_x(d)=\inf_{z\in\R}\bigl(h_K(x,z)-zd\bigr).
 \label{eq:slice-conjugacy}
\end{equation}
Indeed $h_K(x,z)=\sup_d(k_x(d)+zd)$; the reverse formula is concave
biconjugacy, or equivalently separation of the closed hypograph of $k_x$.

We will also retain the following information about optimizers: if
$\xi_m=(x_m,y_m)\to(x,y)$ and $d_m$ maximizes the transform at
$\lambda_0+a_m\xi_m$, then every limit $d$ of the $d_m$ satisfies
\begin{equation}
 d\in\pi_2K,
 \qquad
 r(x,y)=k_x(d)+yd+\frac{\beta_0}{2}d^2.
 \label{eq:scaled-optimizer}
\end{equation}

Write $\beta_m=\beta_0+a_mx_m$ and $h_m=h_0+a_my_m$.
The physical uniform bound on optimizing densities makes $(d_m)$
bounded. For every fixed $z\in\R$, Fenchel's inequality at
$(\beta_m,h_0+a_mz)$ gives
\begin{equation}
 B_m(\xi_m)\leq q_m(x_m,z)+(y_m-z)d_m
                     +\frac{\beta_m}{2}d_m^2.
 \label{eq:scaled-upper}
\end{equation}
This will yield the upper bound after taking a density limit and then
the infimum over $z$.

For the lower bound, choose $v_*=(s_*,d_*)\in K$ that attains the
maximum defining $r(x,y)$, and put $z=y+\beta_0d_*$. Comparing its value
with that at any $(s',d')\in K$ and completing the square gives
\begin{equation}
 x(s_*-s')+z(d_*-d')
       \geq\frac{\beta_0}{2}(d'-d_*)^2.
 \label{eq:strict-face-support}
\end{equation}
Thus $v_*\in F_{(x,z)}$, and every member of this face has density $d_*$.
Choose a full subgradient of $Q_m$ at $(\beta_m,h_0+a_mz)$, and denote
its density coordinate by $d'_m$. Step~1 shows that $d'_m\to d_*$;
nondecreasing dependence on $h$ gives $d'_m\geq0$. Fenchel equality at
$d'_m$ permits evaluation of the transformed objective and yields
\begin{align}
 B_m(\xi_m)
 &\geq q_m(x_m,z)+(y_m-z)d'_m
                       +\frac{\beta_m}{2}(d'_m)^2 \notag\\
 &\longrightarrow xs_*+yd_*+\frac{\beta_0}{2}d_*^2
   =r(x,y).
 \label{eq:scaled-lower}
\end{align}

Now pass to a subsequence in~\eqref{eq:scaled-upper} along which
$d_m\to d$. Taking the infimum over $z$ after the limit and using
\eqref{eq:slice-conjugacy} gives
\[
 \limsup_m B_m(\xi_m)
   \leq k_x(d)+yd+\frac{\beta_0}{2}d^2
   \leq r(x,y).
\]
The finite lower bound~\eqref{eq:scaled-lower} excludes
$d\notin\pi_2K$. Equality throughout proves
\eqref{eq:scaled-optimizer}. These arguments apply to every convergent
sequence $\xi_m$ in a fixed compact set. They therefore prove the uniform
convergence~\eqref{eq:scaled-transform-limit} as well.

We have identified the exact effect of an attraction of strength $a_m$
on the pressure at distance $O(a_m)$ from $\lambda_0$. The positive
quadratic in~\eqref{eq:scaled-transform-limit-definition} is now visible
at order one. We next use it to locate a nearby subdifferential containing
$g$ and to keep all of its reachable gradients within the three groups.

\smallskip
\noindent\textit{Step 3: a nearby point supporting $g$.}
By~\eqref{eq:persistence-input} and $\beta_0>0$,
\begin{equation}
 r(\xi)-g\cdot\xi
    \geq h_K(\xi)-g\cdot\xi
    \geq\gamma|\xi|.
 \label{eq:scaled-coercivity}
\end{equation}
Choose $S>0$ such that $\gamma S>r(0)+1$.
Let $\xi_m$ minimize $B_m(\xi)-g\cdot\xi$ on the closed ball
$\overline B(0,S)$. Uniform convergence on this ball and
\eqref{eq:scaled-coercivity} show that, for large $m$, all such minimizers
are interior: every boundary value exceeds the value at $0$. Every
limit of these minimizers is also interior. Hence, at
\[
 \lambda'_m=\lambda_0+a_m\xi_m,
\]
convex minimality gives
\begin{equation}
 g\in\partial(\T_{a_m}Q_m)(\lambda'_m),
 \qquad |\lambda'_m-\lambda_0|\leq Sa_m.
 \label{eq:nearby-support}
\end{equation}
For large $m$ these points lie in $\operatorname{int}W$.

\smallskip
\noindent\textit{Step 4: localization of every reachable gradient.}
Put $G_m=G_{\T_{a_m}Q_m}(\lambda'_m)$. We claim that, eventually,
\begin{equation}
 G_m\subset\bigcup_{i=1}^3
                  \left(K_i+\overline B(0,\zeta/2)\right).
 \label{eq:reachable-localization}
\end{equation}
Suppose otherwise, and choose $u_m\in G_m$ outside this union along a
subsequence. By the definition of a reachable gradient, there is a
differentiability point $\widehat\lambda_m$ of $\T_{a_m}Q_m$ such that
\[
 |\widehat\lambda_m-\lambda'_m|<a_m^2,
 \qquad
 |\nabla(\T_{a_m}Q_m)(\widehat\lambda_m)-u_m|<1/m.
\]
Write $\widehat\lambda_m=\lambda_0+a_m\widehat\xi_m$, with
$\widehat\xi_m=(\widehat x_m,\widehat y_m)$.
The scaled points remain in a fixed compact set. Moreover, scaling
preserves the gradients of the transforms:
\[
 \nabla B_m(\widehat\xi_m)
       =\nabla(\T_{a_m}Q_m)(\widehat\lambda_m).
\]
Their boundedness thus follows from the local uniform convex convergence
in~\eqref{eq:scaled-transform-limit}. After passing to a subsequence,
assume
\[
 \xi_m\to\xi=(x,y),\qquad
 \widehat\xi_m\to\xi,\qquad u_m\to u=(s,d).
\]

Let $d_m$ be any maximizing density at
$\widehat\lambda_m=(\widehat\beta_m,\widehat h_m)$. Apply
Lemma~\ref{lem:transform-support} with $Q=Q_m$ and $a=a_m$.
Because $H_{Q_m}(\,\cdot\,,d_m)$ is finite and concave, we can choose
$e_m\in-\partial_\beta^+H_{Q_m}(\widehat\beta_m,d_m)$; the lemma gives
\begin{equation}
 \nabla(\T_{a_m}Q_m)(\widehat\lambda_m)
       =\left(e_m+\frac{a_m}{2}d_m^2,d_m\right),
 \label{eq:transformed-gradient}
\end{equation}
and
\begin{equation}
 (e_m,d_m)\in
 \partial Q_m\bigl(\widehat\beta_m,
                   \widehat h_m+\widehat\beta_m a_md_m\bigr).
 \label{eq:full-gradient-support}
\end{equation}
In particular $d_m\to d$, and the first-coordinate increment tends to zero.
The argument on the right has scaled coordinates
\[
 \left(\widehat x_m,
       \widehat y_m+\widehat\beta_m d_m\right)
   \longrightarrow\eta=(x,y+\beta_0d).
\]
Equations~\eqref{eq:transformed-gradient} and
\eqref{eq:full-gradient-support}, together with Step~1, show that
$u\in F_\eta$.

By~\eqref{eq:scaled-optimizer}, there is a point
$v_*=(s_*,d)\in K$ that maximizes the expression defining $r(x,y)$.
Equation~\eqref{eq:strict-face-support} shows that $v_*\in F_\eta$ as
well. The vectors $u$ and $v_*$ have the same density and the same value
against $\eta$, so $xs=xs_*$. Consequently $u$ itself maximizes that
expression. This conclusion also holds when $x=0$; an exposed face need
not be a singleton.

Since $u$ is a maximizer, the geometric observation at the start of the proof
puts $u$ in one of the $K_i$. This contradicts the choice of $u_m$ outside
the fixed $\zeta/2$ neighborhoods. This proves
\eqref{eq:reachable-localization}.

\smallskip
\noindent\textit{Step 5: the three groups and a uniform threshold.}
Define
\[
 J_{i,m}=\conv\left(G_m\cap
                 \left(K_i+\overline B(0,\zeta/2)\right)\right).
\]
Reachable-gradient sets are compact. Hence every nonempty $J_{i,m}$ is
compact and convex, and lies in $K_i+\overline B(0,\zeta/2)$.
The reachable-gradient identity and
\eqref{eq:reachable-localization} give
\[
 \partial(\T_{a_m}Q_m)(\lambda'_m)
       =\conv(J_{1,m}\cup J_{2,m}\cup J_{3,m}).
\]
None of the three sets is empty: otherwise
\eqref{eq:nearby-support} would put $g$ in the convex hull of two of the
allowed enlarged groups. The density coordinate of every reachable
gradient is nonnegative because $\T_{a_m}Q_m$ is nondecreasing in $h$.
Thus the $J_{i,m}$ have all the asserted properties.

We have proved the conclusion for all sufficiently large indices along
every sequence $a_m\downarrow0$, with the same radius $S$. If no uniform
threshold $a_*$ existed, one could choose a sequence of counterexamples
with $a_m\downarrow0$, contradicting what was just proved. Decrease
$a_*$ if necessary so that $a_*\leq\min(1,A+1)$ and
$\lambda_0+\overline B(0,Sa_*)\subset\operatorname{int}W$.
This completes the proof.
\end{proof}

The lemma preserves the three groups even when they contain curved
families of supporting vectors. Its use in the construction is now
precise: an approximation with error at most $a^2$ can be followed by an
attraction of mass $a$, while moving the supporting point by at most
$Sa$ and enlarging each group by an arbitrarily prescribed amount.
